% Full pages of running text, mostly Regular with some italic, bold and small
% caps, set as a book would be: no headings, no paragraph breaks.
%   lualatex --output-directory=build/proofs proofs/textpages.tex   (or make proof)
\documentclass[10.5pt]{article}
\usepackage[paperwidth=148mm,paperheight=210mm,textwidth=104mm,textheight=166mm,
  top=19mm,inner=20mm]{geometry}
\usepackage[french,ngerman,english]{babel}
\usepackage{fontspec}
\usepackage{microtype}
\setmainfont{MyiaDidot-Regular.otf}[Path=../fonts/otf/, Renderer=HarfBuzz, Numbers=OldStyle,
  BoldFont=MyiaDidot-Bold.otf, ItalicFont=MyiaDidot-Italic.otf, BoldItalicFont=MyiaDidot-BoldItalic.otf]
\newcommand\smc[1]{{\addfontfeature{Letters=SmallCaps}#1}}
\newcommand\allsc[1]{{\addfontfeature{Letters={SmallCaps,UppercaseSmallCaps}}#1}}
\linespread{1.12}
\setlength{\parindent}{0pt}
\pagestyle{empty}
\frenchspacing
\begin{document}
\smc{When the Didot family} set up their presses in Paris in the last decades
of the eighteenth century, the printed page in France had already been through
two hundred years of quiet change, and almost nobody outside the trade noticed
what happened next. François-Ambroise Didot, the father, was a printer before he
was a designer of type, and his first concern was not the shape of the letters
but the way they were measured: the \emph{point Didot}, which he fixed at a
little more than the English point, gave the founders a common unit, so that a
line of type cast in one shop would stand true against a line cast in another.
It sounds like a small thing. It was not. A system of sizes that everyone could
rely on meant that a book could be set from the founts of three different
foundries without the compositor having to pack every line with slips of card,
and it meant, too, that the type itself could be made finer, because nothing had
to be thickened to survive a careless fit. The son, Firmin, took the next step.
His roman, cut in the 1780s and refined for twenty years afterwards, pushed the
contrast between thick and thin strokes further than anyone had dared: the
hairlines are almost threads, the serifs flat and unbracketed, the stress
upright, so that a capital \emph{O} looks less like a written circle than like
a drawing made with a ruling pen. Readers who were used to the warmer letters
of Garamond or Caslon found it cold, and some said so in print. Others found it
exactly right for an age that admired order and clear outlines, and the great
editions of Virgil, Horace and Racine that came from the Didot press between 1798
and 1801 were bought, admired and imitated across Europe. In Parma, Giambattista
Bodoni was working along the same lines at the same time, and for a while the
two houses watched each other with a mixture of rivalry and respect; Bodoni
had the more theatrical sense of display, Didot the steadier hand for text. It
is worth remembering how much of what we now call the \textbf{modern face} was
decided by the paper and the press rather than by taste alone. Hairlines that
thin can only be printed on smooth, hard paper, with a well-made press and good
ink, and it was the arrival of wove paper, smoother than the laid sheets of the
earlier century, that made them possible at all. On rough paper the same
letters break up; set too small, they glitter and tire the eye. The designers
knew this. Firmin Didot cut separate sizes for separate uses, each with its own
proportions, and a fount of eight point was not simply a reduced copy of the
fount of twenty-four: its hairlines were heavier, its letters wider, its
spacing more generous. A digital revival has to make the same choices with
other means, and anybody who has set a page of a Didot at the size of a
footnote, on an office printer, has seen what happens when it does not. The
italic presents a different problem. In the eighteenth century the italic was
still, at heart, a separate letter, descended from the chancery hands of the
Renaissance, and it kept its own rhythm long after it had been made to share
the roman’s height and weight. A Didot italic is narrow, steeply sloped and
written rather than drawn; its \emph{a} and \emph{e} are single movements of
the pen, its \emph{f} descends below the line, and its \emph{k} and \emph{v}
end in small curls that have no counterpart in the roman. Setting a phrase in
italic in the middle of a roman line is therefore not a matter of tilting the
letters, and the best compositors used it sparingly: for the titles of books
such as \emph{Paul et Virginie} or \emph{The Vicar of Wakefield}, for words
borrowed from another language, \emph{à propos}, \emph{tout court},
\emph{Zeitgeist}, and for the occasional word that needed a stress the
sentence could not give it. Small capitals had another job. They were used for
the first words of a chapter, for names in a dramatic text, for the
abbreviations that would otherwise shout from the middle of a line, such as
\allsc{UNESCO}, \allsc{NATO} or \allsc{BBC}, and for dates like 1450~\allsc{AD}
or 753~\allsc{BC}. A good fount of small caps is not a set of reduced
capitals: it is drawn to the weight of the lower case, a little wider than the
capitals and a little taller than the x-height, so that a word like
\smc{Paris} or \smc{Amsterdam} sits in the line without a hole and without a
blot. Bold type, finally, is a latecomer, and it was not made for books at all.
It grew out of the jobbing trade of the nineteenth century, out of posters,
handbills and the lists of prices in shop windows, where a single word had to
be seen from across the street, and it took a long time to find its way into
the text page. In a dictionary (the \textbf{\emph{Dictionnaire de l’Académie}}, say) or a reference work it is indispensable: the
entry \textbf{quire}, for instance, is easier to find on a crowded page when
it stands out in bold, and so is the number of a law, \textbf{§~12}, or the
heading word of a footnote. In continuous reading it should be rare. A page
that is full of bold words looks like a page that is shouting, and the reader
soon stops listening. The figures raise questions of their own. Old-style
figures, with their ascenders and descenders, belong in running text, where
they behave like lower-case letters: the year 1789, the price of a quarto
volume at 36~francs, a print run of 2\,500 copies, a page number like p.~147.
Lining figures, all the height of the capitals, belong in tables and in lines
of capitals, where a 1789 that rose and fell would look out of place. The
first editions of the Didot classics were set in numbers that hardly anyone
could afford: the folio Virgil of 1798 was printed in 250 copies, and a copy on
vellum cost more than a working family earned in a year. Yet the letters
spread far beyond those books. By the middle of the nineteenth century, almost
every newspaper in France, Spain and Italy was set in a face that descended
from Didot or Bodoni, and the same was true of a good part of the Greek and
Russian printing of the time. The Russian reader of 1850 met the modern face in
the novels of Gogol, \emph{Мёртвые души} among them, and in the journals of
Петербург; the Greek reader met it in
the editions of the classics produced in Leipzig and Paris, in letters that
followed the Didot Greek closely. Those letters have their own story, and it
is not a simple one. Firmin Didot’s Greek was cut to sit beside his roman, and
it gave up most of the loops and ligatures of the older Greek founts, though
the separate letters we now expect replaced them only slowly, a little at a
time. When a Greek word appears in an English sentence today, it is usually
set in a type that has made peace between the two traditions: λόγος, ἦθος,
\emph{καιρός}. The same is true of the Cyrillic
alphabet, whose civil type, introduced by Peter the Great in 1708, was itself
modelled on the Dutch romans of its time, and which followed the Latin letters
through every change of fashion that came afterwards. A Russian Didot is
therefore not an exotic thing. It is a cousin, cut with the same tools, for
readers who looked at the same engravings and the same architecture, and who
lived, for the most part, in the same idea of what a clear page should be. That
idea has not lost its force. Whatever we think of the coldness of the modern
face, the page it produced was the first that was designed as a whole, from
the measure of the type to the width of the margins, and the habits of order
that it taught are still with us every time we open a well-made book, a good
newspaper or a careful report. The modern face did not survive the
twentieth century unchanged. Photocomposition made every size from a single
drawing, and the hairlines that had been cut heavier for the small sizes were
now reduced with the rest, so that the Didots of the 1970s, admired on the
covers of fashion magazines, were all but unreadable in the footnotes of a
book. The digital fonts that followed inherited the same habit, and for a long
time a designer who wanted a Didot for text had to choose between a face made
for display and one so softened that it had lost the very contrast that gave
it its character. More recent revivals have gone back to the sources. They
measure the punches and the printed pages in the collections of the
\allsc{BNF} in Paris and of the Museo Bodoniano in Parma, they compare the
sizes against one another, and they ask what a letter cut for eight point
would look like if it were redrawn today, with the precision of a screen but
for the softness of paper. The answers are rarely the same twice. One
designer keeps the flat serifs and opens the counters; another lowers the
contrast a little and lengthens the descenders; a third adds an optical size
for captions and another for headlines, and leaves the reader to decide which
belongs where. What they share is a sense that the face deserves to be read,
and not only looked at. A book set in such a type should not draw attention
to itself. On the first page the reader may notice the sharpness of the
letters, the long thin strokes of the \emph{y} and the round, heavy dots on
the \emph{i} and the \emph{j}; by the tenth page, if the work has been done
well, he or she will notice nothing at all, and will simply go on reading, as
the readers of 1800 did, without stopping to wonder why the page seems so
calm. That is the only test that matters in the end, and it is the hardest to
pass, because nobody can see it from a single specimen, a single word, or a
single beautiful capital. It can only be seen in the long run, page after
page, in the middle of a chapter, at the hour when the light is poor and the
reader is tired, and the letters have to carry the sense on their own. Nothing else, no ornament
and no clever arrangement, can do that work for them.
\newpage
\begin{otherlanguage}{french}
\smc{Il faut bien le dire}, la typographie française du début du
dix-neuvième siècle doit presque tout à une seule famille. Les Didot étaient à
la fois imprimeurs, libraires, fondeurs et fabricants de papier, et c’est sans
doute cette réunion des métiers qui leur a permis d’aller si loin : quand on
fabrique soi-même le papier sur lequel on imprimera, on peut graver des
déliés que personne d’autre n’oserait. Firmin Didot ne s’est pas contenté de
dessiner des caractères ; il a perfectionné la stéréotypie, ce procédé qui
consiste à couler une plaque de métal d’après une page composée, et qui permit
de tirer à bas prix les classiques que réclamait un public nouveau. On lui doit
aussi, avec son père, le \emph{point typographique} qui porte encore leur nom
et qui a servi, pendant près de deux siècles, à mesurer les corps dans toute
l’Europe continentale. Les éditions dites \emph{du Louvre}, parce qu’elles
furent imprimées dans les galeries du palais où le gouvernement avait installé
l’atelier, sont restées célèbres : le Virgile, l’Horace et le Racine en trois
volumes, tirés à quelques centaines d’exemplaires, furent salués, à
l’une des premières expositions des produits de l’industrie, comme les plus
beaux livres jamais imprimés. Le
jugement paraît aujourd’hui excessif, mais il dit bien l’ambition de
l’entreprise. On y voit une page où rien n’est laissé au hasard : les marges
sont vastes, les capitales espacées avec soin, l’italique réservé aux titres
et aux citations, les petites capitales employées pour les noms des
personnages, et l’on y chercherait en vain un caractère gras, que l’on
considérait alors comme un procédé d’affiche. Il est d’ailleurs frappant de
voir combien la lettre de Didot s’accorde avec la langue française : les
accents aigus, graves et circonflexes, les cédilles et les trémas, si nombreux
dans une page de Racine, trouvent leur place au-dessus de lettres dont le
dessin est assez régulier pour les recevoir sans désordre. Qu’on lise à haute
voix une phrase comme « où est-il allé, ce garçon qu’on croyait déjà à
Noël ? », et l’on verra que l’œil, lui aussi, a besoin de ces signes pour ne
pas trébucher. C’est dans cet accord, plus encore que dans la finesse des
déliés, qu’il faut chercher la raison d’une si longue fortune : deux siècles plus
tard, le lecteur d’un roman de poche retrouve, sans le savoir, la page que
l’on composait chez les Didot.
\end{otherlanguage}
\newpage
\begin{otherlanguage}{ngerman}
\smc{Im deutschen Sprachraum} hatte es die Antiqua der Didot zunächst schwer.
Bis weit ins neunzehnte Jahrhundert hinein wurden die meisten Bücher in
Fraktur gesetzt, und die runde, helle Schrift aus Paris galt vielen Lesern als
fremd, ja als Zeichen einer Gesinnung, die man nicht teilen wollte. Wer sie
dennoch verwendete, tat es mit Absicht. Der Leipziger Drucker und Verleger
Georg Joachim Göschen etwa ließ die Werke Wielands in einer Antiqua drucken,
die sich unmittelbar an Didot anlehnte, und schuf damit einige der schönsten
Bücher seiner Zeit; Goethe, der sich für solche Fragen lebhaft interessierte,
verfolgte die Arbeit mit großer Aufmerksamkeit. Für wissenschaftliche Werke,
für Übersetzungen aus dem Lateinischen und Griechischen und für alles, was
sich an ein gelehrtes Publikum in ganz Europa richtete, setzte sich die
Antiqua nach und nach durch, während die Fraktur für Zeitungen, Gesangbücher
und Volksausgaben noch lange die Regel blieb. Diese Spaltung hat Spuren in der
Typografie hinterlassen, die bis heute sichtbar sind. Das lange \emph{s} etwa,
das in der Fraktur selbstverständlich war, wurde in der Antiqua schon früh
aufgegeben; das scharfe \emph{ß} dagegen blieb, und mit ihm die Frage, wie es
in einer Zeile aus Großbuchstaben oder Kapitälchen aussehen soll. Wer heute
eine Überschrift wie \allsc{STRASSE} oder \allsc{GRÖSSE} setzt, steht noch vor
derselben Wahl, die die Setzer des neunzehnten Jahrhunderts beschäftigte, auch
wenn es inzwischen ein großes \emph{ẞ} gibt, das seit 2017 offiziell zugelassen
ist. Auch die Umlaute stellten die Schriftschneider vor Aufgaben: Über einem
\emph{ä}, \emph{ö} oder \emph{ü} müssen zwei Punkte Platz finden, ohne an die
Zeile darüber zu stoßen, und bei den Großbuchstaben \emph{Ä}, \emph{Ö} und
\emph{Ü} ist der Raum noch knapper. Ein guter Satz nimmt darauf Rücksicht,
ohne dass der Leser es bemerkt, und genau darin liegt seine Kunst. Dass die
Antiqua sich am Ende durchgesetzt hat, lag weniger an ästhetischen Gründen als
an der Geschichte selbst: Nach 1941 wurde die Fraktur in Deutschland von
Amts wegen abgeschafft, und nach dem Krieg wollte niemand sie zurück. Die
Schriften der Didot und Bodoni aber, die man so lange für fremd gehalten
hatte, waren inzwischen längst zu Hause, in Zeitungen, Büchern und
Briefköpfen, und kaum ein Leser wusste noch, woher sie gekommen waren.
\end{otherlanguage}
\end{document}
